\documentclass{article}
\usepackage[utf8]{inputenc}
\usepackage{amsmath}
\usepackage{geometry}
\geometry{a4paper, margin=1in}
\usepackage{parskip}
\usepackage{enumitem}
\usepackage{hyperref}
\usepackage[T1]{fontenc}
\usepackage{lmodern}

\begin{document}

\title{Movie Recommender System with Guided Active Learning}
\author{Project Team}
\date{July 15, 2025}
\maketitle

\section{Introduction}
This project develops a guided active learning setup for cold-start movie recommendations using the MovieLens small dataset, which contains 100,836 ratings of 9,742 movies by 610 users. The system learns user preferences through interactive ratings while providing feedback on recommendation impacts, aiming to enhance user engagement and recommendation accuracy.

\section{Task 1: Guided Active Learning Setup}
\subsection{Methodology}
The recommender system uses matrix factorization to model user-movie interactions:
\[
\min_{U,V} \sum_{(i,j) \in K} (R_{ij} - U_i^T V_j)^2 + \lambda (\|U\|_F^2 + \|V\|_F^2)
\]
where $R_{ij}$ is the rating by user $i$ for movie $j$, $U_i \in \mathbb{R}^K$ and $V_j \in \mathbb{R}^K$ are latent representations, and $\lambda = 0.1$ is the regularization parameter. The latent dimension is set to $K = 10$.

For a new user, the system learns their latent representation $U_i$ by solving:
\[
\min_{U_i} \sum_{j \in K_i} (R_{ij} - U_i^T V_j)^2 + \lambda \|U_i\|_F^2
\]
Movies are selected for rating using uncertainty sampling, choosing those with the highest prediction variance based on the current user vector. In the experimental condition, users receive feedback showing the top 5 recommended movies based on their ratings, updated dynamically.

\subsection{Implementation}
The system is implemented as a Django app within the `project4` module, using Python for backend logic in `recommender.py`, HTML templates (`index.html`, `study_interface.html`), and JavaScript (`app.js`) for client-side interactivity. The MovieLens dataset is processed using Pandas, and Math.js handles matrix operations in the browser.

\section{Task 2: User Study Design}
\subsection{Hypothesis}
Providing feedback on how ratings influence recommendations increases user engagement (measured by the number of ratings provided) and improves recommendation accuracy (measured by RMSE on a test set).

\subsection{Design}
\begin{itemize}
    \item \textbf{Participants}: 50 participants recruited via online platforms (e.g., Prolific).
    \item \textbf{Conditions}:
    \begin{itemize}
        \item \textit{Control}: Standard active learning without feedback.
        \item \textit{Experimental}: Active learning with feedback on recommendation impact.
    \end{itemize}
    \item \textbf{Procedure}:
    \begin{enumerate}
        \item Participants complete a consent form and demographic survey.
        \item Rate 10 movies selected by uncertainty sampling.
        \item Experimental group sees dynamic recommendation feedback.
        \item Rate 5 additional movies of their choice (test set).
        \item Complete a post-study survey on engagement (Likert scale).
    \end{enumerate}
    \item \textbf{Metrics}:
    \begin{itemize}
        \item Engagement: Number of ratings provided.
        \item Accuracy: RMSE on test set ratings.
    \end{itemize}
    \item \textbf{Analysis}: Compare metrics between groups using t-tests.
\end{itemize}

\subsection{Recruitment}
Participants will be recruited via Prolific, targeting adults who watch movies regularly. Incentives will be provided per platform guidelines. Inclusion criteria: age 18+, fluent in English.

\section{Task 3: User Study Interface}
The interface, implemented in `study_interface.html`, allows users to select and rate movies. The experimental condition displays a list of top recommended movies after each rating. The interface logs ratings for analysis and displays final recommendations after 10 ratings. The system is served via the Django app, with the MovieLens dataset in `project4/data/`.

\section{Conclusion}
This project implements a novel guided active learning approach for cold-start movie recommendations, with a user study to evaluate its effectiveness. The interface and report are accessible via the landing page (`index.html`).

\bibliography{references}
\bibliographystyle{plain}
\begin{thebibliography}{1}
\bibitem{harper2015}
F. Maxwell Harper and Joseph A. Konstan.
The MovieLens Datasets: History and Context.
\textit{ACM Transactions on Interactive Intelligent Systems (TiiS)}, 5(4):19:1–19:19, 2015.
\url{https://doi.org/10.1145/2827872}.
\end{thebibliography}

\end{document}
